\section{Why tiny-mass padding is insufficient}
\label{sec:padding}

Take $\varphi(x)=x^2$, initially with four locations
\[
 (A,B,C,D)=(0,2,3,5)
\]
and equal, unnormalized masses $1$.  Normalizing all masses multiplies
every cost by the same factor and changes no price.  Here
\[
 q=D(BC)=\tfrac12,\qquad
 Q=D(AB)+D(CD)=4,\qquad
 D(ABC)=D(BCD)=\tfrac{14}{3}.
\]
The four-state deterministic price is exactly $7/6$.  If the fine pair is
$BC$, its two triple coarsenings have cost $14/3$, and its remaining
coarsening $AD\mid BC$ has cost $13$.  Any other fine pair has cost at
least $2$, giving a fine-level ratio at least $4$.  This covers all exact
hierarchies, including noncontiguous ones; hierarchies with unused capacity
can be refined to exact capacities without increasing cost.

Now add $E=-1$ with mass $\varepsilon\in(0,1)$, keeping the four original
masses equal to $1$, and normalize by $4+\varepsilon$.  At capacities three
and four, both flat optima are attained by a single hierarchy:
\[
 H_4=EA\mid B\mid C\mid D,
 \qquad H_3=EA\mid BC\mid D.
\]
Its unnormalized costs are
\[
 q_5=\frac{\varepsilon}{1+\varepsilon}<\frac12,
 \qquad
 Q_5=\frac{\varepsilon}{1+\varepsilon}+\frac12<1.
\]
The pair $EA$ is strictly cheapest: other pairs touching $E$ have the same
mass coefficient and strictly larger squared distance, while all core
pairs cost at least $1/2$.

It remains to check the coarse optimum.  A three-cell partition of five
states has one triple or two pairs.  Every core triple costs at least
$14/3$.  A triple touching $E$ and a core pair other than $BC$ dominates
that core pair and therefore costs at least $2$.  The remaining triple is
$EBC$, of cost
\[
 D(EBC)=\frac{25\varepsilon+1}{\varepsilon+2},\qquad
 D(EBC)-Q_5
 =\frac{\varepsilon(47\varepsilon+45)}
 {2(\varepsilon+1)(\varepsilon+2)}>0.
\]
For two-pair partitions, if one pair is $BC$, the cheapest disjoint pair
is $EA$.  If neither pair is $BC$, at least one pair is a core pair of cost
at least $2$: two disjoint pairs cannot both touch $E$.  These candidates
also cost more than $Q_5$.  Thus the claimed coarse partition is optimal,
and
\[
 \rho_{5,\mathrm{det}}=\rho_{5,\mathrm{stoch}}=1
 \quad\hbox{for every }0<\varepsilon<1.
\]
All five locations can be moved into $[1/4,3/4]$ by
$y=(x+4)/12$.  Positive scaling makes the quadratic curvature at most one,
and the usual curvature-integral construction gives strictly proper losses
in $[0,1/2]$.  Thus neither distinctness, interior support, bounded loss,
nor arbitrarily small padding mass rescues this failed embedding method.


This refutes preservation under naive tiny-mass padding. It does not
contradict the carefully separated anchors in Theorem~\ref{thm:embedding}.
